\documentclass[11pt,a4paper]{article}
 \usepackage[spanish]{babel} 
 \usepackage[utf8]{inputenc} 
 \usepackage[T1]{fontenc}
  \usepackage{amsmath,amssymb,amsthm}
   \usepackage{geometry} 
   \geometry{margin=2.5cm}
   
 \usepackage{tikz}
\usepackage{pgfplots}

\decimalpoint


    \title{Comportamiento dinámico de algunas funciones racionales en el plano complejo} 
    \author{}
     \date{} 
     
     
     \begin{document}
      \maketitle
      
       \section{Comportamiento dinámico de la función \(f(z)=z^2\)} 
       
       La dinámica compleja estudia el comportamiento de las sucesivas iteraciones de una función holomorfa. En este documento analizamos uno de los ejemplos más sencillos y fundamentales: 
       \[ f:\mathbb{C}\longrightarrow\mathbb{C}, \qquad f(z)=z^2. \] 
       
       Dado un punto inicial \(z_0\in\mathbb{C}\), consideramos su órbita hacia delante 
       \[ z_0,\quad z_1=f(z_0),\quad z_2=f(z_1),\quad \ldots, \] 
       es decir, $z_n=f^{n}(z_0),$ donde \(f^{n}\) denota la \(n\)-ésima iterada de \(f\). Nuestro objetivo es describir el comportamiento asintótico de las órbitas. 
       
   La expresión de las iteradas de \(f\) puede obtenerse explícitamente. En efecto, 
   \[ f^2(z)=f(f(z))=(z^2)^2=z^4, \]
    \[ f^3(z)=z^8, \] 
    y, por inducción, \[ f^n(z)=z^{2^n}. \] 
  Por tanto, el estudio dinámico de \(f\) se reduce al análisis de la sucesión $z_n=z^{2^n}. $
       
También es sencillo obtener la representación polar de las mismas.  Sea $z=re^{i\theta},$ donde \(r=|z|\ge 0\). Entonces $f(z) = r^2 e^{2i\theta}$ y, en general, $f^{n}(z) = r^{2^n}e^{\,i2^n\theta}.$ Por tanto, la iteración produce dos efectos: 
\begin{itemize}
 \item el módulo se eleva al cuadrado; 
 \item el argumento se duplica.
  \end{itemize} 

  
 Los únicos puntos fijos finitos  son $z=0$ (superatractor) y $ z=1$ (repulsor). Además el punto del infinito también es un punto fijo (superatractor).
 
 
 
 Si $|z|<1$, se tiene que $f^n(z)\longrightarrow 0$, es decir, todos los puntos del disco unitario abierto convergen al origen. Por contra, si   $|z|>1$, se tiene que $f^n(z)\longrightarrow \infty$, y todos los puntos exteriores al círculo unidad escapan al infinito. 
 
 El comportamiento dinámico más interesante se da en la circunferencia unidad, cuando $|z|=~1$. Si $z=e^{i\theta},$ entonces \[ f^n(z)=e^{i2^n\theta}. \] El módulo permanece constante e igual a \(1\). La dinámica depende del valor de \(\theta\): 
 \begin{itemize} 
 \item Si \(\theta/(2\pi)\) es racional, la órbita es periódica o eventualmente periódica. 
 \item Si \(\theta/(2\pi)\) es irracional, la órbita no es periódica y se distribuye densamente sobre la circunferencia unidad. 
 \end{itemize} 
 
 En consecuencia, en este caso, el conjunto de Fatou, \(F(f)\), consta de dos componentes: 
 \begin{enumerate} 
 \item la cuenca de atracción del origen: el interior del disco unidad.
 \item la cuenca de atracción del infinito: el exterior del disco unidad.
 \end{enumerate} 
 
 El conjunto de Julia \(J(f)\) es la frontera del conjunto de Fatou, en este caso es exactamente la circunferencia unidad
  \[ J(f)=\{z\in\mathbb C:\ |z|=1\}. \] 
  
  Vamos a analizar gráficamente el comportamiento de  las órbitas de 4 puntos iniciales $z_0$ muy próximos entre sí, distinguiendo los casos de que no estén en la circunferencia unidad (figura~\ref{Fig:1}), o que estén en la circunferencia unidad (figura~\ref{Fig:2}).
  
\begin{figure}[htbp]
\begin{center}
\begin{tikzpicture}[scale=2.75]
%% Parámetro c
%\def\cre{0.5}
%\def\cim{0.5}

% Ejes
\draw(-1.1,0) -- (1.1,0);
\draw  (0,-1.1) -- (0,1.1);

%% Ejes
%\draw[->] (-0.6,0) -- (0.6,0) node[right] {$x$};
%\draw[->] (0,-0.6) -- (0,0.6) node[above] {$y$};

%% Cardioide parametrizada
%\draw[thick,blue,smooth,domain=0:360,samples=400]
%  plot ({0.5*cos(\x) - 0.25*cos(2*\x)},
%        {0.5*sin(\x) - 0.25*sin(2*\x)});


%Circunferencia unidad
\draw[thick,blue,smooth,domain=0:360,samples=400]
  plot ({cos(\x)}, {sin(\x)});



% Dependiendo lo que queramos refeljar, podemos poner punto inicial z0=0 o c0=c
\pgfmathsetmacro{\x}{0.86}
\pgfmathsetmacro{\y}{0.5}
%\pgfmathsetmacro{\x}{0}
%\pgfmathsetmacro{\y}{0}

% Dibujar punto inicial
 \fill[red](\x,\y) circle (0.02) node[above right] {$\x+\y i$};

% Iteraciones
\foreach \n in {1,...,20} {

    % Calcular z_{n+1}
    \pgfmathsetmacro{\xn}{\x*\x - \y*\y }
    \pgfmathsetmacro{\yn}{2*\x*\y }

    % Dibujar segmento
    \draw[gray] (\x,\y) -- (\xn,\yn);

    % Dibujar puntos
    \fill[black] (\xn,\yn) circle (0.02);

    % Actualizar valores
    \xdef\x{\xn}
    \xdef\y{\yn}
}
\end{tikzpicture}
\quad
\begin{tikzpicture}[scale=2.75]
%% Parámetro c
%\def\cre{0.5}
%\def\cim{0.5}

% Ejes
\draw(-1.1,0) -- (1.1,0);
\draw  (0,-1.1) -- (0,1.1);

%% Ejes
%\draw[->] (-0.6,0) -- (0.6,0) node[right] {$x$};
%\draw[->] (0,-0.6) -- (0,0.6) node[above] {$y$};

%% Cardioide parametrizada
%\draw[thick,blue,smooth,domain=0:360,samples=400]
%  plot ({0.5*cos(\x) - 0.25*cos(2*\x)},
%        {0.5*sin(\x) - 0.25*sin(2*\x)});


%Circunferencia unidad
\draw[thick,blue,smooth,domain=0:360,samples=400]
  plot ({cos(\x)}, {sin(\x)});



% Dependiendo lo que queramos refeljar, podemos poner punto inicial z0=0 o c0=c
\pgfmathsetmacro{\x}{0.87}
\pgfmathsetmacro{\y}{0.495}
%\pgfmathsetmacro{\x}{0}
%\pgfmathsetmacro{\y}{0}

% Dibujar punto inicial
 \fill[red](\x,\y) circle (0.02) node[below] {$\x+\y i$};

% Iteraciones
\foreach \n in {1,...,9} {

    % Calcular z_{n+1}
    \pgfmathsetmacro{\xn}{\x*\x - \y*\y }
    \pgfmathsetmacro{\yn}{2*\x*\y }

    % Dibujar segmento
    \draw[gray] (\x,\y) -- (\xn,\yn);

    % Dibujar puntos
    \fill[black] (\xn,\yn) circle (0.02);

    % Actualizar valores
    \xdef\x{\xn}
    \xdef\y{\yn}
}
\end{tikzpicture}
\caption{Dos instantáneas de las órbitas de $z^2$: a la izquierda, órbita de $z_0=0.86+0.5i$, punto situado en el interior de la circunferencia unidad. A la derecha, órbita de $z_0=0.87+0.495i$, punto situado en el exterior de la circunferencia unidad. Se aprecia la convergencia al origen en el primer caso, mientras que en el segundo caso, la órbita escapa al infinito. }
\label{Fig:1}
\end{center}
\end{figure}
  
  BLA, BLA, BLA

BLA, BLA, BLA

BLA, BLA, BLA

BLA, BLA, BLA

BLA, BLA, BLA

BLA, BLA, BLA

BLA, BLA, BLA

BLA, BLA, BLA

BLA, BLA, BLA

BLA, BLA, BLA

BLA, BLA, BLA



  \begin{figure}[htbp]
\begin{center}
\begin{tikzpicture}[scale=2.75]
%% Parámetro c
%\def\cre{0.5}
%\def\cim{0.5}

% Ejes
\draw(-1.1,0) -- (1.1,0);
\draw  (0,-1.1) -- (0,1.1);

%% Ejes
%\draw[->] (-0.6,0) -- (0.6,0) node[right] {$x$};
%\draw[->] (0,-0.6) -- (0,0.6) node[above] {$y$};

%% Cardioide parametrizada
%\draw[thick,blue,smooth,domain=0:360,samples=400]
%  plot ({0.5*cos(\x) - 0.25*cos(2*\x)},
%        {0.5*sin(\x) - 0.25*sin(2*\x)});


%Circunferencia unidad
\draw[thick,blue,smooth,domain=0:360,samples=400]
  plot ({cos(\x)}, {sin(\x)});



% Dependiendo lo que queramos reflejar, podemos poner punto inicial z0=0 o c0=c
%\pgfmathsetmacro{\x}{sqrt(3)/2}
%\pgfmathsetmacro{\y}{1/2}
\pgfmathsetmacro{\x}{cos(30)} %Los ángulos van en grados
\pgfmathsetmacro{\y}{sin(30)}

% Dibujar punto inicial. Hago trampa y lo pongo a mano
 \fill[red](\x,\y) circle (0.02) node[above right] {$e^{\pi i/6}$};%{$\x+\y i$};

% Iteraciones
\foreach \n in {1,...,20} {

    % Calcular z_{n+1}
    \pgfmathsetmacro{\xn}{\x*\x - \y*\y }
    \pgfmathsetmacro{\yn}{2*\x*\y }

    % Dibujar segmento
    \draw[gray] (\x,\y) -- (\xn,\yn);

    % Dibujar puntos
    \fill[black] (\xn,\yn) circle (0.02);

    % Actualizar valores
    \xdef\x{\xn}
    \xdef\y{\yn}
}
\end{tikzpicture}
\quad
\begin{tikzpicture}[scale=2.75]
%% Parámetro c
%\def\cre{0.5}
%\def\cim{0.5}

% Ejes
\draw(-1.1,0) -- (1.1,0);
\draw  (0,-1.1) -- (0,1.1);

%% Ejes
%\draw[->] (-0.6,0) -- (0.6,0) node[right] {$x$};
%\draw[->] (0,-0.6) -- (0,0.6) node[above] {$y$};

%% Cardioide parametrizada
%\draw[thick,blue,smooth,domain=0:360,samples=400]
%  plot ({0.5*cos(\x) - 0.25*cos(2*\x)},
%        {0.5*sin(\x) - 0.25*sin(2*\x)});


%Circunferencia unidad
\draw[thick,blue,smooth,domain=0:360,samples=400]
  plot ({cos(\x)}, {sin(\x)});



% Dependiendo lo que queramos refeljar, podemos poner punto inicial z0=0 o c0=c
\pgfmathsetmacro{\x}{0.8664783803811612}
\pgfmathsetmacro{\y}{0.4992145994780599}
%\pgfmathsetmacro{\x}{cos(360*sqrt(2)/17)}
%\pgfmathsetmacro{\y}{sin(360*sqrt(2)/17)}

% Dibujar punto inicial. Lo pongo a mano
 \fill[red](\x,\y) circle (0.02) node[above right] {$e^{(2\sqrt{2}\pi i)/17}$};%{$\x+\y i$};

% Iteraciones
\foreach \n in {1,...,14} {

    % Calcular z_{n+1}
    \pgfmathsetmacro{\xn}{\x*\x - \y*\y }
    \pgfmathsetmacro{\yn}{2*\x*\y }

    % Dibujar segmento
    \draw[gray] (\x,\y) -- (\xn,\yn);

    % Dibujar puntos
    \fill[black] (\xn,\yn) circle (0.02);

    % Actualizar valores
    \xdef\x{\xn}
    \xdef\y{\yn}
}
\end{tikzpicture}
\caption{Dos instantáneas de las órbitas de $z^2$. A la izquierda, órbita de $z_0=e^{\pi i/6}=(\sqrt{3}+ i)/2\approx 0.866025+0.5i$, punto situado en la circunferencia unidad y con una órbita eventualmente periódica (acaba en el 2-ciclo  $\{e^{2\pi i/3},e^{-2\pi i/3}\}$). A la derecha, órbita $z_0=e^{2\sqrt{2}\pi i/17}\approx 0.866478 + 0.499215i$, punto situado en la circunferencia unidad y con una órbita no periódica. }
\label{Fig:2}
\end{center}
\end{figure}

NOTA: en el segundo dibujo hacemos trampas. Al  trabajar con una aproximación de $e^{2\sqrt{2}\pi i/17}$, llega un momento en que las órbitas abandonan la circunferencia unidad, por eso solo hemos cogido 14 iteraciones


BLA, BLA, BLA

BLA, BLA, BLA

BLA, BLA, BLA

BLA, BLA, BLA

BLA, BLA, BLA

BLA, BLA, BLA

BLA, BLA, BLA

BLA, BLA, BLA

BLA, BLA, BLA

BLA, BLA, BLA

BLA, BLA, BLA

BLA, BLA, BLA

BLA, BLA, BLA

BLA, BLA, BLA

BLA, BLA, BLA

BLA, BLA, BLA


En la figura~{Fig:3} se muestra un ejemplo para incluir gráficas importadas.


\begin{figure}[bht]
\centering 
\rule{0.45\textwidth}{0.45\textwidth}% %%%% <--- Borrador sin dibujos
\qquad
\rule{0.45\textwidth}{0.45\textwidth}% %%%% <--- Borrador sin dibujos
%\includegraphics[width=0.45\textwidth]{figuratal.pdf}% %%%% <--- Con dibujos
%\qquad
%\includegraphics[width=0.45\textwidth]{figuracual.pdf}% %%%% <--- Con dibujos
\caption{Ejemplo para incluir dos figuras, figuratal.pdf y figuracual.pdf.}
\label{Fig:3}
\end{figure}



BLA, BLA, BLA

BLA, BLA, BLA

BLA, BLA, BLA

BLA, BLA, BLA

BLA, BLA, BLA

BLA, BLA, BLA

BLA, BLA, BLA

BLA, BLA, BLA

BLA, BLA, BLA

BLA, BLA, BLA


BLA, BLA, BLA

BLA, BLA, BLA

BLA, BLA, BLA

BLA, BLA, BLA

BLA, BLA, BLA

BLA, BLA, BLA

BLA, BLA, BLA

BLA, BLA, BLA

BLA, BLA, BLA

BLA, BLA, BLA

\begin{figure}[p] 
\centering 
\setlength{\tabcolsep}{3pt} 
\begin{tabular}{ccc} 
\includegraphics[width=0.31\textwidth]{Pella.png} & \includegraphics[width=0.31\textwidth]{Pella.png} & \includegraphics[width=0.31\textwidth]{Pella.png} \\[4mm] \includegraphics[width=0.31\textwidth]{Pella.png} & \includegraphics[width=0.31\textwidth]{Pella.png} & \includegraphics[width=0.31\textwidth]{Pella.png} \\[4mm] \includegraphics[width=0.31\textwidth]{Pella.png} & \includegraphics[width=0.31\textwidth]{Pella.png} & \includegraphics[width=0.31\textwidth]{Pella.png} \\[4mm] 
%\includegraphics[width=0.31\textwidth]{Pella.png} & \includegraphics[width=0.31\textwidth]{Pella.png} & \includegraphics[width=0.31\textwidth]{Pella.png} 
\end{tabular} 
\caption{Conjunto de 9 figuras en una distribución $3\times 3$, para que se muestren en una página. Poner [p] después de begin figure} 
\label{fig:nueve-figuras} 
\end{figure}

\begin{figure}[p] 
\centering 
\setlength{\tabcolsep}{3pt} 
\begin{tabular}{ccc} 
\includegraphics[height=0.2\textheight]{Pella.png} & \includegraphics[height=0.2\textheight]{Pella.png} & \includegraphics[height=0.2\textheight]{Pella.png} \\[4mm] \includegraphics[height=0.2\textheight]{Pella.png} & \includegraphics[height=0.2\textheight]{Pella.png} & \includegraphics[height=0.2\textheight]{Pella.png} \\[4mm] 
\includegraphics[height=0.2\textheight]{Pella.png} & \includegraphics[height=0.2\textheight]{Pella.png} & \includegraphics[height=0.2\textheight]{Pella.png} \\[4mm] 
\includegraphics[height=0.2\textheight]{Pella.png} & \includegraphics[height=0.2\textheight]{Pella.png} & \includegraphics[height=0.2\textheight]{Pella.png} 
\end{tabular} 
\caption{Conjunto de 12 figuras en una distribución $4\times 3$, para que se muestren en una página. Poner [p] después de begin figure} 
\label{fig:doce-figuras} 
\end{figure}


\end{document}

